% =====================================================================
% Section: Experiments -- Comparison with state-of-the-art (outline.md §3.2;
%   defines sec:main_comparison). Discusses Table I (tab:main_comparison) and
%   references Table II (tab:efficiency).
% Draft v1 (2026-06-17). Every numeric statement is read VERBATIM from
%   paper/tables/table1_main_comparison.tex (self-tested DPRNet rows; competitor
%   rows cited from original papers per the table notes). Bold/underline status in
%   the table is the basis for "best/second-best" wording.
% HONEST FRAMING (matches corrected data): against the strongest baseline CATANet,
%   DPRNet attains the best result on nearly every benchmark/scale cell, slightly
%   surpassing CATANet (margins ~0.01-0.07 dB) while CATANet is consistently second.
%   The ONLY metric where CATANet still leads is x4 Manga109 SSIM (0.9183 vs 0.9181).
%   Margins are small -- frame as "matches or slightly surpasses", not a large gain.
%   DPRNet also carries a modestly higher parameter count than CATANet (660K vs 535K).
% =====================================================================
\subsection{Comparison with state-of-the-art}
\label{sec:main_comparison}

We compare DPRNet with eight representative lightweight SR methods. The CNN group
contains CARN~\cite{ahn2018carn}, IMDN~\cite{hui2019imdn}, RFDN~\cite{liu2020rfdn},
and RLFN~\cite{kong2022rlfn}, while the Transformer and routing group contains
SwinIR-light~\cite{liang2021swinir}, ELAN-light~\cite{zhang2022elan},
SRFormer-light~\cite{zhou2023srformer}, and the direct predecessor
CATANet~\cite{liu2025catanet}. Competitor numbers are cited from the original
papers, as noted in Table~\ref{tab:main_comparison}. DPRNet uses one unified
checkpoint per scale and is reported on all five benchmarks without per-dataset
selection.

\paragraph{Overall accuracy.}
Table~\ref{tab:main_comparison} shows that DPRNet and CATANet form the top tier at
all three scales, both clearly ahead of the remaining seven methods. For example,
at $\times4$ DPRNet reaches $32.59$/$0.9002$\,dB on Set5 and $27.76$/$0.7434$ on
BSD100, exceeding the strongest non-routing competitor SRFormer-light
($32.51$/$0.8988$ and $27.73$/$0.7422$) while using a comparable parameter budget.
The consistent advantage of the two content-routing methods over window-based and
distillation-based baselines supports the value of content-aware token
aggregation in the lightweight regime.

\paragraph{Comparison with CATANet.}
Against CATANet---the strongest baseline and the method from which DPRNet differs
only in the routing module---DPRNet attains the best result on nearly every
benchmark/scale cell, with CATANet consistently ranked second. The margins are small but
consistent: at $\times4$ DPRNet improves Urban100 from $26.87$/$0.8081$ to
$26.89$/$0.8089$, Set5 from $32.58$/$0.8998$ to $32.59$/$0.9002$, Set14 from
$28.90$/$0.7880$ to $28.91$/$0.7883$, and Manga109 PSNR from $31.31$ to $31.32$;
at $\times2$ it reaches the best Set14 ($34.06$/$0.9220$), Urban100 PSNR ($33.15$),
and Manga109 ($39.38$/$0.9789$); and at $\times3$ it leads or ties CATANet on every
benchmark (e.g.\ BSD100 $29.29$/$0.8103$ and Urban100 $29.06$/$0.8689$). The single
metric on which CATANet remains ahead is Manga109 SSIM at $\times4$
($0.9183$ vs.\ $0.9181$). These gaps are within a few hundredths of a dB. The
result therefore supports a conservative conclusion: replacing the EMA-center routing of TAB with the dynamic
prototype routing of DPR matches and slightly surpasses CATANet-level
reconstruction quality while changing the routing mechanism to the interpretable,
confidence-aware formulation analyzed in Sec.~\ref{sec:ablation} and
Sec.~\ref{sec:routing_analysis}.

\paragraph{Efficiency.}
Table~\ref{tab:efficiency} reports the efficiency profile at $\times4$. DPRNet has
$660$K parameters and $52.14$\,G Multi-Adds for a fixed $256{\times}256$ LR input,
which keeps it within the lightweight bracket alongside the compared methods. Under
the same input convention, it uses fewer Multi-Adds than the lightweight Transformer
baselines SwinIR-light ($60.3$\,G) and SRFormer-light ($56.5$\,G), while delivering top-tier
accuracy. Both its parameter count and Multi-Adds are modestly higher than CATANet's
($660$K/$52.14$\,G vs.\ $535$K/$46.8$\,G at $\times4$), reflecting the learnable query
bank and router projections of DPR, so no efficiency advantage over CATANet is claimed. The
per-scale profile is given in Table~\ref{tab:efficiency_profile}. Since several
competitor Multi-Adds and all competitor latencies are not available under a shared
measurement protocol, those entries are intentionally left unreported rather than mixed
across incompatible hardware or input conventions.
